\section{有关工作}

\paragraph{无模型 Atari}

应用于 Atari 基准的大多数智能体都使用无模型算法训练。DQN \citep{mnih2015dqn} 表明，在 Q-learning 中引入经验回放和目标网络后，可以训练深度神经网络策略。随后多项工作扩展了 DQN，例如 DDQN 中的偏差校正\citep{van2015ddqn}、优先经验回放\citep{schaul2015prioritizedreplay}、架构改进\citep{wang2016duelingdqn}，以及分布式价值学习\citep{bellemare2017c51,dabney2017qrdqn,dabney2018iqn}。除价值学习外，基于策略梯度的智能体也被用于 Atari 基准，例如 ACER\citep{schulman2017ppo}、PPO\citep{schulman2017ppo}、ACKTR\citep{wu2017acktr} 和 Reactor\citep{gruslys2017reactor}。另一类工作通过分布式数据收集提升性能，并常常把环境步预算提高到 2 亿以上\citep{mnih2016a3c,schulman2017proximal,horgan2018apex,kapturowski2018r2d2,badia2020agent57}。

\paragraph{世界模型}

若干基于模型的智能体关注本体感知输入\citep{watter2015e2c,gal2016deeppilco,higuera2018synthesizing,henaff2018planbybackprop,chua2018pets,tingwu2019benchmark,wang2019poplin}；也有工作对图像建模，但不把该模型用于规划\citep{oh2015atari,krishnan2015deepkalman,karl2016dvbf,chiappa2017recenvsim,babaeizadeh2017sv2p,gemici2017temporalmemory,denton2018stochastic,buesing2018dssm,doerr2018prssm,gregor2018tdvae}；还有工作结合基于模型和无模型方法的优势\citep{kalweit2017blending,nagabandi2017mbmf,weber2017i2a,kurutach2018modeltrpo,buckman2018steve,ha2018worldmodels,wayne2018merlin,igl2018dvrl,srinivas2018upn,lee2019slac}。\citet{risi2019discrete} 使用进化搜索优化离散潜变量；\citet{parmas2019pipps} 结合了 Reinforce 梯度和重参数化梯度。到目前为止，带图像输入的大多数世界模型智能体仍局限于相对简单的控制任务\citep{watter2015e2c,ebert2017visualmpc,ha2018worldmodels,hafner2018planet,zhang2018solar,hafner2019dreamer}。下文详细介绍两种已应用于 Atari 的基于模型方法。

\begin{table}[t!]
\vspace*{-3ex}
\small
\centering
\begin{tabular}{@{}lccccccc@{}}
\toprule
\textbf{算法} &
\makecell[c]{\textbf{奖励}\\\textbf{模型制作}} &
\makecell[c]{\textbf{图像}\\\textbf{模型制作}} &
\makecell[c]{\textbf{潜在}\\\textbf{转移}} &
\makecell[c]{\textbf{单}\\\textbf{GPU}} &
\makecell[c]{\textbf{可训练}\\\textbf{参数}} &
\makecell[c]{\textbf{Atari}\\\textbf{帧数}} &
\makecell[c]{\textbf{加速器}\\\textbf{天数}} \\
\midrule
DreamerV2 & \cmark & \cmark & \cmark & \cmark & 22M & 200M & 10 \\
SimPLe & \cmark & \cmark & \xmark & \cmark & 74M & 4M & 40 \\
MuZero & \cmark & \xmark & \cmark & \xmark & 40M & 20B & 80 \\
MuZero Rean\rlap{alyze} & \cmark & \xmark & \cmark & \xmark & 40M & 200M & 80 \\
\bottomrule
\end{tabular}
\caption{近期利用学习模型进行规划的 RL 算法概念比较。DreamerV2 和 SimPLe 利用图像输入提供的学习信号来学习完整环境模型，而 MuZero 通过单个任务特有的价值梯度学习其模型。MuZero 使用的蒙特卡洛树搜索很有效，但增加了复杂性，也难以并行化。该组件与本文提出的世界模型是正交的。}
\label{tab:properties}
\end{table}


\paragraph{SimPLe}

SimPLe 智能体\citep{kaiser2019simple} 在像素空间学习视频预测模型，并用其预测来训练 PPO 智能体\citep{schulman2017ppo}，如 \cref{tab:properties} 所示。该模型根据前四帧直接预测下一帧，并接收一个额外离散潜变量作为输入。作者在 Atari 子集上以 40 万和 200 万环境步评估 SimPLe，之后报告收益递减。近期一些无模型方法也沿用了 40 万步设置进行比较\citep{srinivas2020curl,kostrikov2020drq}。然而，在这种数据效率设置中达到的最高性能只是玩家归一化中位数 0.28\citep{kostrikov2020drq}，距离人类水平仍很远。相比之下，我们关注 2 亿帧后的成熟且竞争性更强的评估，此时有许多成功无模型算法可供比较。

\paragraph{MuZero}

MuZero 智能体\citep{schrittwieser2019muzero} 学习奖励和值的序列模型\citep{oh2017vpn}，并通过 Monte Carlo 树搜索解决强化学习任务\citep[MCTS;][]{coulom2006efficient,silver2017alphago}。如 \cref{tab:properties} 所示，该序列模型完全通过预测任务相关信息训练，并不包含利用图像的显式表示学习。MuZero 表明，在大量工程投入和庞大计算预算下，规划可以在若干棋盘游戏和确定性 Atari 游戏中取得强性能。不过，MuZero 并未公开实现，而且在单个 GPU 上训练一个 Atari 智能体需要两个多月。相比之下，DreamerV2 是一个较简单的算法，能在单个 GPU 上 10 天内达到 Atari 人类水平，因此更容易被许多研究者复现。此外，MuZero 的高级规划组件与 DreamerV2 学得的准确世界模型互补，未来可以结合。DreamerV2 还利用输入图像提供的额外学习信号，这类似于近期半监督图像分类中的成功经验\citep{chen2020simclr,he2020moco,grill2020byol}。
